\documentclass[11pt]{article}
\usepackage{fontspec,amsmath,amssymb,amsthm,hyperref}
\usepackage[margin=1in]{geometry}
\setlength{\emergencystretch}{3em}
\title{Dalalyan’s theorem for the actual Gaussian ULA transition}
\author{LangevinLean mathematical exposition}
\date{}
\begin{document}\maketitle
\noindent Source module: \texttt{UnadjustedLangevin/CanonicalULA.lean}. Compile this proof document with XeLaTeX or LuaLaTeX. The accompanying HTML places the complete Lean declarations beside every proof.


Let \(U:\mathbb R^d\to\mathbb R\) be differentiable and \(m\)-strongly convex with \(m>0\), with \(L\)-Lipschitz gradient. The canonical Brownian source derives all stochastic reference obligations and proves the Wasserstein bound for the iterated Gaussian transition \(P_h\), with finite initial \(W_2\).


\section{Identify the canonical output law}


On each initial-extension product, Brownian source innovations have independent standard normal scalar coordinates, hence iid standard Gaussian vector laws, independently of initialization. The lifted-chain law theorem therefore identifies the common product output with \(P_h^k\mu\), where \(P_h\mu=\operatorname{Law}(X-h\nabla U(X)+\sqrt{2h}Z)\). Finite coupling and Gibbs-moment premises justify its \(L^2\) realization.



\noindent Lean declarations: \texttt{productBrownianOutput\_eq\_ulaLawStep}.

\section{Derive stationarity for every reference segment}


Construct a global integral solution on a horizon containing the desired segment. BrownianSegmentsGlobal identifies the chosen segment with a time shift of that solution. GibbsStationarity proves that each global value has Gibbs law, using the Metropolis refinement argument and product independence from the reference initial state. Thus independent Brownian coordinates build every ProductReference field, including stationarity, from the potential assumptions. The explicit canonical Brownian construction supplies one such source, and its normalized increments define canonicalEulerNoise.



\noindent Lean declarations: \texttt{productReference\_of\_independent\_brownian}, \texttt{canonicalEulerNoise}, \texttt{canonicalProductReference}.

\section{Guard finiteness and expose the final bound}


A finite initial Wasserstein distance yields a finite initial coupling. The canonical-output identity rewrites the product output as \(P_h^k\mu\); the product-reference finiteness theorem proves \(W_2(P_h^k\mu,\pi)<\infty\). In positive dimension, strong convexity and Lipschitzness imply \(m\le L\). Insert the canonical reference into the product theorem to obtain, for \(0<h\le2/(m+L)\),

\[W_2(P_h^k\mu,\pi)\le(1-mh)^kW_2(\mu,\pi)+\frac{33}{20}\frac Lm\sqrt{hd}.\]

For \(2/(m+L)\le h<2/L\), the same assembly gives

\[W_2(P_h^k\mu,\pi)\le(Lh-1)^kW_2(\mu,\pi)+\frac{33}{20}\frac{Lh}{2-Lh}\sqrt{hd}.\]

All real-valued distances in these statements have separately established finiteness, so extended-real conversion cannot hide an infinite transport cost.



\noindent Lean declarations: \texttt{canonicalOutput\_eq\_ulaLawStep}, \texttt{dalalyan\_output\_finite}, \texttt{dalalyan}, \texttt{dalalyan\_large\_step}.

\end{document}
